\documentclass[11pt]{article}
\usepackage[T1]{fontenc}
\usepackage{textcomp}
\usepackage[margin=0.75in]{geometry}
\usepackage{amsmath,amssymb,amsthm}
\usepackage{array,booktabs}
\usepackage{hyperref}
\setlength{\parindent}{0pt}
\newtheorem{theorem}{Theorem}
\theoremstyle{definition}
\newtheorem{definition}{Definition}
\title{Flow Proof Artifact: List-len-append-singleton}
\author{Generated by \texttt{flow doc proof}}
\date{Flow Proof Book}
\begin{document}
\maketitle
Appending a singleton increases length by one.\\[0.5em]
\textbf{Source.} church — https://en.wikipedia.org/wiki/Length\_of\_a\_list\\[0.5em]
\bigskip
\noindent{\large \textbf{Derived fact 1.} \textit{len(xs ++ cons(y, nil)) = succ(len(xs))} \hfill List \cdot length \cdot append singleton increases length by one}\par
\smallskip\noindent\textit{Coordinate: List \cdot length \cdot append singleton increases length by one}\par
\smallskip\noindent\textit{Source: church}\par
\smallskip\noindent\textit{Built on: length adds under append, for length on List, successor on the right via commutativity, for addition on the natural numbers}\par
\medskip
\noindent\textbf{Goal.} len(xs ++ cons(y, nil)) = succ(len(xs))\par
\noindent$\displaystyle \forall xs \in List \forall y \in \mathbb{N}\quad len(xs ++ cons(y, nil)) = \mathrm{succ}(len(xs))$\par
\medskip
\noindent
\begin{tabular}{@{} >{\raggedright\arraybackslash}p{0.06\textwidth} >{\raggedright\arraybackslash}p{0.47\textwidth} >{\raggedleft\arraybackslash}p{0.41\textwidth} @{}}
\textbf{\#} & \textbf{Proof} & \textbf{Mathematics} \\
\midrule
\textbf{1.} & We prove that append singleton increases length by one for length on List. & \\[0.45em]
\textbf{2.} & We invoke the derived fact governing length on List: length adds under append, for length on List (instantiated for xs, cons(y, nil)). & \\[0.45em]
\textbf{3.} & We invoke the derived fact governing addition on the natural numbers: successor on the right via commutativity, for addition on the natural numbers (instantiated for len(xs), 0). & \\[0.45em]
\textbf{4.} & From step 2 and step 3, this implies len(xs plus plus cons(y, nil)) equals the successor of len(xs). Hence proven. & $\displaystyle len(xs ++ cons(y, nil)) = \mathrm{succ}(len(xs))$ \\[0.45em]
\bottomrule
\end{tabular}
\label{thm:List-length-append_singleton_increases_length_by_one}
\par\medskip\hrule
\end{document}
